\documentclass[11pt,a4paper]{article}

\usepackage[utf8]{inputenc}
\usepackage[T1]{fontenc}
\usepackage[french]{babel}
\usepackage{geometry}
\usepackage{booktabs}
\usepackage{hyperref}

\geometry{margin=2.5cm}
\setlength{\parskip}{0.5em}
\setlength{\parindent}{0pt}

\title{TP 3 --- Clustering de documents textuels}
\author{Hany El Atlassi}
\date{\today}

\begin{document}

\maketitle

\section{Démarche et résultats}

J'ai appliqué un clustering non supervisé (K-Means) sur le corpus BBC News
fourni dans \texttt{data/HANY EL ATLASSI - Travail a faire 2 Kmeans .xlsx},
pour voir si un algorithme n'ayant jamais vu les étiquettes des documents
parvient à retrouver leurs 5 catégories d'origine : \texttt{sport} (511
articles), \texttt{business} (510), \texttt{politics} (417), \texttt{tech}
(401) et \texttt{entertainment} (386), pour un total de 2225 articles. Le
code complet se trouve dans \texttt{main.ipynb}.

J'ai d'abord transformé chaque article en vecteur TF-IDF (mots vides
anglais retirés, termes trop rares ou trop fréquents filtrés), ce qui m'a
donné une matrice de 2225 documents pour 9136 termes. J'ai fixé le nombre
de clusters à $k=5$, un choix que j'ai confirmé avec la courbe du coude et
le score de silhouette, et qui correspond au nombre de catégories réelles
afin de pouvoir comparer directement les deux. J'ai ensuite entraîné
K-Means (\texttt{k-means++}, 10 initialisations) sur cette matrice.

En regardant les 10 termes les plus proches de chaque centroïde, j'ai pu
identifier le thème dominant de chaque cluster :

\begin{table}[h]
\centering
\small
\begin{tabular}{cll}
\toprule
Cluster & Termes principaux & Thème \\
\midrule
0 & people, mobile, technology, users, music, digital, games, phone, software, net & tech \\
1 & game, england, win, cup, match, team, players, injury, play, club & sport \\
2 & film, best, awards, award, band, festival, music, actor, star, album & entertainment \\
3 & mr, labour, election, blair, party, brown, government, howard, minister, tory & politics \\
4 & mr, growth, government, company, economy, market, bank, sales, oil, firm & business \\
\bottomrule
\end{tabular}
\end{table}

J'ai ensuite croisé ces clusters avec les catégories réelles, ce qui donne
le tableau de contingence suivant :

\begin{table}[h]
\centering
\small
\begin{tabular}{lrrrrr}
\toprule
Catégorie \textbackslash{} Cluster & 0 & 1 & 2 & 3 & 4 \\
\midrule
business      &   7 &   0 &   0 &   3 & 500 \\
entertainment &   5 &   2 & 321 &   0 &  58 \\
politics      &   4 &   1 &   0 & 274 & 138 \\
sport         &   0 & 479 &   2 &   0 &  30 \\
tech          & 363 &   8 &   5 &   0 &  25 \\
\bottomrule
\end{tabular}
\end{table}

J'ai constaté que chaque cluster a bien un unique thème dominant, et j'ai
calculé deux indices de comparaison : un Adjusted Rand Index (ARI) de 0.695
et une Normalized Mutual Information (NMI) de 0.743.

\section{Analyse et interprétation}

Le clustering n'a accès qu'à la proximité des vecteurs TF-IDF : un document
est placé dans le cluster dont le centroïde partage le plus de termes
fréquents et discriminants avec lui. J'ai remarqué que les catégories dont
le vocabulaire est très spécifique se retrouvent naturellement bien
isolées : le sport a son propre lexique (\emph{match, team, club, win}) qui
n'apparaît presque jamais ailleurs, ce qui explique pourquoi j'ai retrouvé
479 des 511 articles \texttt{sport} (94\%) dans le même cluster. À
l'inverse, j'ai observé que \texttt{business} et \texttt{politics}
partagent un vocabulaire institutionnel et économique commun (\emph{mr,
government, growth, market}), car la presse traite souvent les mêmes
événements (budget, décisions gouvernementales, croissance) sous ces deux
angles à la fois. C'est ce chevauchement lexical, et non un défaut de
l'algorithme, qui explique pourquoi 138 articles \texttt{politics} se
retrouvent absorbés par le cluster \texttt{business}, et pourquoi 58
articles \texttt{entertainment} (souvent liés à l'industrie du
divertissement en tant qu'activité économique) le sont également.

J'ai aussi voulu vérifier si le déséquilibre des classes de départ (511
articles pour \texttt{sport} contre 386 pour \texttt{entertainment}, un
rapport d'environ 1.3) expliquait les erreurs observées, mais ce n'est pas
ce que j'ai constaté : la catégorie la plus mal reconstituée
(\texttt{business}, avec 500 documents correctement isolés sur 510 mais 268
documents d'autres classes qui lui sont rattachés) n'est ni la plus grande
ni la plus petite. J'en conclus que c'est surtout la proximité sémantique
entre thèmes voisins qui pèse sur le résultat ; un déséquilibre plus marqué
aurait en revanche pu favoriser les grandes classes, K-Means cherchant à
minimiser une inertie globale que les clusters les plus peuplés dominent
davantage.

J'ai choisi l'ARI et la NMI parce que ce sont deux indices externes conçus
pour comparer un clustering à un étiquetage de référence sans supposer que
les identifiants de cluster correspondent aux identifiants de classe (rien
n'oblige le cluster 4 à s'appeler \texttt{business}). L'ARI corrige en plus
l'accord attendu par simple hasard, ce qui m'a semblé plus fiable qu'un
taux de concordance brut, surtout que mes classes n'ont pas toutes la même
taille. La valeur de 0.695 que j'ai obtenue me montre que l'accord entre
mon clustering et les catégories réelles est nettement supérieur à ce
qu'un partitionnement aléatoire produirait, tout en restant loin de
l'accord parfait, ce qui correspond bien aux confusions que j'ai observées
entre \texttt{business} et \texttt{politics}. La NMI mesure quant à elle la
quantité d'information partagée entre les deux partitions, normalisée
entre 0 et 1 : la valeur de 0.743 que j'ai obtenue m'indique que connaître
le cluster d'un document réduit fortement mon incertitude sur sa catégorie
réelle, ce qui confirme que le vocabulaire TF-IDF capture bien la
structure thématique du corpus, malgré les zones de recouvrement lexical
entre certains sujets proches.

\end{document}
